\scanpage{185}
\mspage{5\deleted{\unsure{8}}9}
% archive sheet number: 184

prenant $M_0=M\cap\mathrm{Filt}^{\unsure{\Pi|\pi}}$. D'autre part, $c'\Rightarrow c$, (\unsure{sans} hyp. sur $\pi$\unsure{?}), car \deleted{\unsure{donc}} \unsure{OPS|OBS} que la partie $M_0$ de $\mathrm{Filt}^{\Pi}$ est stable par $\underline{\lim}$ filtrantes de cardinal $\leq\pi$. Mais alors, connaissant l'équivalence de a,b,c, d'où \inserted{\unsure{aussitôt|aussi}} le corollaire, il \deleted{est} s'ensuit que l'enveloppe stable par $\underline{\lim}$ filtrantes de $M_0$ dans $\mathcal M$ \deleted{est} n'est autre que la partie accessible \deleted{\unsure{en}} déduite de $M_0$ dont il est question dans c, et on gagne.

Donc \deleted{\unsure{tout}} revient à prouver que, si tous les objets de $\mathrm{Filt}^{\Pi}$ sont accessibles, alors c implique a). Montrons d'abord que c) implique le critère d'appartenance \inserted{à $M$} (5.3.) (p. 51.), donc par la \deleted{démonstration \unsure{précédente} de} prop. 5.6. ``revue et corrigée'' en marge, il s'ensuit déjà que c) $\Rightarrow$ b), et il restera à établir b) $\Rightarrow$ a) i.e. que b) implique la stabilité de $M$ par $\underline{\lim}$ filtrantes quelconques (pas \unsure{seult} \unsure{gdes|grdes} devant $\pi$).

Que c) implique 5.3. va résulter du lemme suivant, appliqué à $I=\mathrm{Ss}_\pi(X)$ :

\uline{Lemme 5.8.} Soit $I$ un ensemble ordonné \uline{$\pi$-adapté}, i.e. filtrant et tel que toute partie filtrante $K$ de $I$, \inserted{de cardinal $\leq\pi$,} ait une borne supérieure $i_K$ dans $I$. \deleted{Considérons} \inserted{Je dis} \deleted{\unsure{l'ensemble des}} qu'une partie \deleted{filtrante} $J$ de $I$ est \deleted{strictement} \uline{$\pi$-cofinale} si elle est \deleted{filtrante} \inserted{cofinale} et si pour toute partie filtrante $K$ de $J$ de cardinal $\leq\pi$, la borne supérieure $i_K$ dans $I$ est

